\documentclass[11pt]{article}

\usepackage{amsmath,amssymb,amsthm,mathtools}
\usepackage[margin=1in]{geometry}
\usepackage{url}

\newtheorem{theorem}{Theorem}
\newtheorem{quotedtheorem}{Quoted theorem}
\newtheorem{remark}{Remark}

\newcommand{\R}{\mathbb{R}}
\newcommand{\Vol}{\operatorname{Vol}}
\newcommand{\Dil}{\operatorname{Dil}}

\begin{document}
\title{Submission C solution to Problem 4}
\date{}
\maketitle


\begin{abstract}
We prove the stated estimate as an immediate consequence of Guth's
product-domain form of his area-contraction estimates for rectangles.
All constants below are positive dimensional constants; in this problem
the dimension is fixed equal to \(4\), so the final constant is universal.
\end{abstract}

\section{A product estimate of Guth}

We shall use the following special case of Guth's product-domain estimate.
It is the case \(n=4\), \(k=2\), \(\dim X=1\), and \(\dim Y=2\) of
Proposition 9.1 in Guth's paper \cite[Proposition 9.1]{GuthAreaExpanding}.
The proof in \cite{GuthAreaExpanding} is obtained from the isoperimetric
profile of rectangles, the Almgren--Gromov complex-of-cycles formalism,
and the tightening argument developed in Sections 1, 3, 4, and 6 of that
paper; a fuller account appears in Guth's thesis \cite{GuthThesis}.

\begin{quotedtheorem}[Guth, specialized product estimate]\label{qt:guth}
There are constants \(\varepsilon_0>0\) and \(c_0>0\) with the following
property.

Let
\[
X=[0,a],\qquad I=[0,b],\qquad Y=[0,u]\times[0,v],
\]
with the product Euclidean metric, and let \(S\subset \R^4\) be a rectangle
with side lengths
\[
S_1\leq S_2\leq S_3\leq S_4.
\]
Suppose
\[
F:(X\times I\times Y,\partial(X\times I\times Y))\longrightarrow
(S,\partial S)
\]
is piecewise smooth, has non-zero degree, and satisfies
\[
\Dil_2(F)\leq 1.
\]
If
\[
a\leq \varepsilon_0 S_1
\qquad\text{and}\qquad
uv\leq \varepsilon_0 S_3S_4,
\]
then
\[
\Vol_4(X\times I\times Y)
=abuv
\geq
c_0\,S_1S_2^{1/2}S_3^{3/2}S_4.
\]
\end{quotedtheorem}

\begin{remark}
Guth states his estimates for Lipschitz/PL complexes of cycles. A
piecewise smooth map on a compact piecewise smooth rectangle is admissible
in that framework, since push-forwards of integral Lipschitz chains are
defined simplexwise. Alternatively, one may approximate a piecewise smooth
relative map by smooth relative maps with \(2\)-dilation increased by an
arbitrarily small factor and then let the approximation parameter tend to
zero. The relative degree is unchanged under such approximation.
\end{remark}

\section{Proof of the required estimate}

\begin{theorem}\label{thm:main}
There is a universal constant \(\kappa>0\) such that the following holds.
Let \(R\subset \R^4\) and \(S\subset \R^4\) be rectangles with side lengths
\[
R_1\leq R_2\leq R_3\leq R_4,
\qquad
S_1\leq S_2\leq S_3\leq S_4.
\]
Suppose
\[
f:(R,\partial R)\longrightarrow (S,\partial S)
\]
is piecewise smooth, has degree \(1\), and satisfies \(\Dil_2(f)\leq 1\).
If
\[
R_1\leq \kappa S_1,
\]
then either
\[
R_3R_4>\kappa S_3S_4,
\]
or
\[
R_1R_2R_3R_4
>
\kappa\,S_1S_2^{1/2}S_3^{3/2}S_4.
\]
\end{theorem}

\begin{proof}
Let \(\varepsilon_0,c_0\) be the constants from
Theorem~\ref{qt:guth}. Define
\[
\kappa:=\frac12\min\{\varepsilon_0,c_0\}.
\]
Assume
\[
R_1\leq \kappa S_1.
\]
If
\[
R_3R_4>\kappa S_3S_4,
\]
then the first desired alternative holds, and there is nothing to prove.

Thus assume instead that
\[
R_3R_4\leq \kappa S_3S_4.
\]
Write the domain rectangle as the product
\[
R=[0,R_1]\times[0,R_2]\times\bigl([0,R_3]\times[0,R_4]\bigr).
\]
Set
\[
X=[0,R_1],\qquad I=[0,R_2],\qquad
Y=[0,R_3]\times[0,R_4].
\]
Then
\[
\Vol_1(X)=R_1\leq \kappa S_1\leq \varepsilon_0S_1
\]
and
\[
\Vol_2(Y)=R_3R_4\leq \kappa S_3S_4\leq \varepsilon_0S_3S_4.
\]
The map \(f\) is a degree-one map of pairs
\[
f:(X\times I\times Y,\partial(X\times I\times Y))
\longrightarrow (S,\partial S)
\]
with \(\Dil_2(f)\leq 1\). Hence Theorem~\ref{qt:guth} applies and gives
\[
R_1R_2R_3R_4
=
\Vol_4(X\times I\times Y)
\geq
c_0\,S_1S_2^{1/2}S_3^{3/2}S_4.
\]
Since \(\kappa\leq c_0/2\), this implies
\[
R_1R_2R_3R_4
>
\kappa\,S_1S_2^{1/2}S_3^{3/2}S_4.
\]
Therefore the second desired alternative holds. This proves the theorem.
\end{proof}

\begin{thebibliography}{9}

\bibitem{GuthAreaExpanding}
L.~Guth,
\emph{Area-expanding embeddings of rectangles},
arXiv:0710.0403, 2007.

\bibitem{GuthThesis}
L.~Guth,
\emph{Area-contracting maps between rectangles},
Ph.D. thesis, Massachusetts Institute of Technology, 2005.

\bibitem{FedererFleming}
H.~Federer and W.~H. Fleming,
\emph{Normal and integral currents},
Ann. of Math. (2) \textbf{72} (1960), 458--520.

\bibitem{Almgren}
F.~J. Almgren,
\emph{The homotopy groups of the integral cycle groups},
Topology \textbf{1} (1962), 257--299.

\bibitem{GromovFilling}
M.~Gromov,
\emph{Filling Riemannian manifolds},
J. Differential Geom. \textbf{18} (1983), no.~1, 1--147.

\end{thebibliography}

\end{document}